\documentclass[11pt, a4paper]{article}
\usepackage[utf8]{inputenc}
\usepackage{amsmath}
\usepackage{amssymb}
\usepackage{amsfonts}
\usepackage{graphicx}
\usepackage{booktabs}
\usepackage{multirow}
\usepackage[colorlinks=true, linkcolor=blue, citecolor=blue, urlcolor=blue]{hyperref}
\usepackage{geometry}
\geometry{margin=1in}

\title{\textbf{SpatialIntegrator: A Foundation Model-Driven Framework for Multimodal Integration of Histopathology and Spatial Transcriptomics}}

\author{
    \parbox{\textwidth}{\centering
        \textbf{Ronie Martínez}$^{1,2,*}$\quad
        \textbf{Santiago Guerrero}$^{2}$\quad
        \textbf{Jennyfer García-Cárdenas}$^{2}$\quad
        \textbf{Lizbeth Peña-Zúñiga}$^{2}$ \\[2ex]
        $^{1}$ \small Departamento de Informática y Ciencias de la Computación, \\
        Escuela Politécnica Nacional, Quito, Ecuador \\[1ex]
        $^{2}$ \small Laboratorio de Ciencia de Datos Biomédicos, \\
        Facultad de Ciencias Médicas de la Salud y de la Vida, \\
        Universidad Internacional del Ecuador, Quito, Ecuador \\[1.5ex]
        $^*$ \small Corresponding author: \texttt{ronie.martinez@epn.edu.ec}
    }
}

\date{}

\begin{document}

\maketitle

\begin{abstract}
Spatial transcriptomics (ST) technologies have transformed translational oncology by enabling whole-transcriptome quantification within native histopathologically preserved architecture. Despite capturing paired high-resolution morphological images (e.g., Hematoxylin and Eosin, H\&E) alongside genomic arrays, dominant analytical workflows routinely disregard histological features during tissue domain discovery, relying exclusively on highly sparse mRNA expression count matrices. This unimodal dependence exacerbates technical sequencing dropouts, lateral transcript diffusion artifacts, and microenvironmental blurring across malignant boundaries. To address these limitations, we introduce SpatialIntegrator, an open-source multimodal computational framework that bridges deep visual representations with spatial gene expression through self-supervised pathology Vision Foundation Models (VFMs). By deploying histopathology-specialized architectures—including Phikon and UNI, self-supervisedly pre-trained on millions of digital pathology whole-slide image tiles—SpatialIntegrator extracts dense morphological representations directly at the spot receptive field. We introduce an orthogonal algebraic fusion mechanism incorporating Frobenius norm inertia equalization, which mathematically unifies variance-stabilized transcriptional profiles with visual embeddings via tunable parametric weighting ($\alpha \in [0,1]$) prior to high-speed graph-based Leiden community detection. Comprehensive evaluation on standardized 10x Genomics Visium human infiltrating ductal carcinoma datasets demonstrates that SpatialIntegrator achieves up to a 10-fold enhancement in macro-architectural spatial contiguity compared to unimodal RNA baselines (elevating Spatial Silhouette Scores from $-0.0536$ to $+0.2141$). Furthermore, utilizing rigorous non-parametric Wilcoxon rank-sum testing, morphological prioritization ($\alpha=0.2$, $224\times224$px receptive tile resolution) accurately isolates invasive tumor perimeter niches enriched for hallmark lipogenic and epithelial oncogenes (\textit{ERBB2}, \textit{FASN}) as well as stromal matrix-remodeling enzymes (\textit{MMP11}, \textit{COL1A1}), confirming that histopathological awareness acts as a biological denoising filter without eclipsing transcriptional specificity or artificially penalizing naturally dispersed multifocal microenvironments. Designed for accessible democratization, SpatialIntegrator integrates \texttt{scanpy}, \texttt{squidpy}, and \texttt{timm} backends into an interactive Streamlit platform available at \url{https://github.com/MrOnie/SpatialIntegrator}.
\end{abstract}

\section{Introduction}

The spatial coordinates of heterogeneous cellular phenotypes within solid tumors dictate invasive progression, immune evasion, and therapeutic resistance. Spatially resolved transcriptomics (ST)—recognized as Nature's Method of the Year \cite{marx2021method}—has catalyzed a paradigm shift by empowering quantitative mapping of whole-transcriptome expression across intact tissue architectures \cite{stahl2016visualization, rao2021exploring, moses2022museum}. High-throughput multi-cellular array technologies, such as 10x Genomics Visium and Slide-seq, capture paired experimental modalities: high-resolution optical histological staining (Hematoxylin and Eosin, H\&E) co-registered with barcoded sequencing spot arrays \cite{stahl2016visualization, palla2022squidpy}.

Despite the dual molecular-morphological nature of these datasets, established bioinformatics pipelines treat tissue imagery as an incidental aesthetic background. Routine computational workflows for spatial domain identification and cell-type niche discovery rely exclusively on dimensionality reduction manifolds—such as Principal Component Analysis (PCA) and Uniform Manifold Approximation and Projection (UMAP)—computed solely over the raw transcriptomic count matrix \cite{wolf2018scanpy, hao2021integrated}. However, RNA-centric clustering faces substantial technical vulnerabilities. Array-based ST platforms experience high zero-count dropout rates, uneven sequencing coverage, lateral mRNA diffusion during enzymatic tissue permeabilization, and spatial multiplexing (where individual $55\mu\text{m}$ spots homogenize transcripts from 5 to 15 heterogeneous cells) \cite{larsson2021spatially}. Consequently, unimodal RNA clustering routinely yields fragmented, physiologically inconsistent domain classifications that conflate distinct anatomical microenvironments—such as desmoplastic tumor stroma, necrotic cores, and proliferating ductal parenchyma \cite{andersson2020single}.

Concurrent with advancements in genomics, digital computational pathology has undergone an epochal transition driven by self-supervised Vision Foundation Models (VFMs). Whereas early morphological quantification relied on supervised Convolutional Neural Networks (CNNs) constrained by localized phenotypic annotation schemas \cite{he2016deep, kather2019predicting}, modern architectures leverage Vision Transformers (ViTs) scaled through masked image modeling and contrastive self-supervised distillation \cite{dosovitskiy2021image, he2022masked}. Foundation models including UNI \cite{chen2024general} (pre-trained over 100 million clinical pathology patches) and Phikon \cite{filiot2024scaling} (distilled via iBOT across millions of pan-cancer histology tiles from the Cancer Genome Atlas) encode robust, highly generalizable features of cellular morphology, nuclear pleomorphism, and stromal extracellular matrices without requiring explicit clinical labels.

Recent authoritative consensus across computational biology underscores that deep learning-based multimodal integration of histology and ST represents the cornerstone of next-generation spatial analysis \cite{luo2025deep, zhang2024inferring, yang2025staig}. Pioneering frameworks including stLearn \cite{pham2020stlearn}, SpaGCN \cite{hu2021spagcn}, and STNet \cite{he2020integrating} demonstrated early efficacy by uniting classical CNN image features with gene expression. Most recently (2025--2026), the literature has advanced toward uniting digital pathology foundation models directly with spatial transcriptomics, exemplified by state-of-the-art frameworks including FmH2ST \cite{fmh2st2025}, SEAL \cite{seal2026}, and STORM \cite{storm2026}. However, while these cutting-edge deep learning frameworks significantly improve representation learning, they exhibit profound computational and operational constraints: they largely depend on deep dual-branch graph neural networks, hierarchical transformer pre-training, or intensive contrastive loss fine-tuning schemas that demand immense multi-GPU overhead, obscure parametric interpretation, and lack lightweight plug-and-play compatibility with standard \texttt{AnnData} clinical infrastructures \cite{wang2025stimage, fu2024bidcell}.

To overcome these computational and translational roadblocks, we introduce \textbf{SpatialIntegrator}. Our modular Python framework establishes direct mathematical integration between self-supervised pathology foundation models (\texttt{timm}, Hugging Face) \cite{wightman2019pytorch} and genomic workflows (\texttt{scanpy}, \texttt{squidpy}) \cite{wolf2018scanpy, palla2022squidpy}. By employing a transparent, algebraically orthogonal latent space fusion mechanism reinforced by Frobenius norm spectral inertia equalization, SpatialIntegrator mitigates transcriptional noise and defines topologically coherent microenvironments. Moreover, we deliver an intuitive Streamlit interactive platform that eliminates programming barriers, enabling wet-lab translational oncologists and computational biologists alike to perform high-resolution, image-aware biological discovery.

\section{Materials and Methods}

\subsection{Data Structures and Genomic Preprocessing}
The SpatialIntegrator platform natively builds upon the robust \texttt{AnnData} object architecture \cite{wolf2018scanpy}. For quantitative empirical benchmarking, human breast infiltrating ductal carcinoma spatial transcriptomics arrays and paired optical high-resolution H\&E slide imagery were obtained from official validated repositories using \texttt{squidpy.datasets.visium\_hne\_adata()} \cite{palla2022squidpy}.

To mitigate variance instability and baseline library size discrepancies across capture spots, the raw mRNA counts matrix was subjected to stringent quantitative preprocessing. Total library size normalization was enacted to scale individual spot read depth to a persistent target sum of $10,000$ counts, followed by natural logarithmic transformation ($\log(x+1)$). To filter ubiquitous baseline housekeeping noise and focus analytical resolution on biological heterogeneity, the top $3,000$ highly variable genes (HVGs) were prioritized utilizing Seurat variance dispersion estimation \cite{stuart2019comprehensive}. Finally, expression distributions across selected HVGs were zero-mean centered and scaled to unit variance, producing the standardized transcriptional matrix $\mathbf{R} \in \mathbb{R}^{n \times g}$, where $n$ denotes spatial capture spots and $g=3,000$ variable genes.

\subsection{Vision Transformer Architecture and Patch Feature Extraction}
To harvest phenotypic tissue architecture, we instantiated an extensible \texttt{ImageEmbedder} class interfacing natively with deep learning model hubs (\texttt{timm} and \texttt{transformers}) \cite{wightman2019pytorch}.

\textbf{Spatial Localization and Tiling:} Physical spatial coordinates $(x_i, y_i)$ corresponding to each sequencing spot $i$, originating within the pixel reference frame of the sequencing matrix, were scaled by the Visium calibration parameter \texttt{tissue\_hires\_scalef} to accurately intercept the high-resolution optical H\&E whole-slide image. Around each transformed centroid, square morphological tile fields of dimension $H \times W$ pixels were systematically cropped. To evaluate receptive field dependency across biological architectures, tile extraction was benchmarked across three quantitative scales: $112\times112$, $224\times224$, and $336\times336$ pixels, capturing subcellular density up to macro-stromal connective tissue boundaries.

\textbf{VFM Representation Inference:} Cropped histological tiles were propagated directly through pre-trained Vision Foundation Models. We evaluated an ImageNet general-purpose baseline (ViT-Base/16) \cite{dosovitskiy2021image} alongside clinical histopathology specialists: UNI \cite{chen2024general} and Phikon \cite{filiot2024scaling}. In standard ViT topology, an input tile is partitioned into fixed non-overlapping spatial token patches ($16\times16$px), linearly projected into embedding space, and appended with a learnable classification token (\texttt{[CLS]}). Following iterative self-attention encoding across $L=12$ (Phikon/ViT-B) or $L=24$ (UNI) transformer blocks, the final hidden state representation of the \texttt{[CLS]} token aggregates the comprehensive semantic morphology of the entire physical tile. Formally, extracting this terminal hidden state yielded a dense morphological vector $\mathbf{v}_i \in \mathbb{R}^d$ for each spatial spot $i$ ($d = 768$ for ViT-Base/Phikon; $d = 1024$ for UNI), constructing the global histological matrix $\mathbf{V} \in \mathbb{R}^{n \times d}$.

\subsection{Orthogonal Algebraic Multimodal Fusion and Frobenius Normalization}
Fusing high-dimensional genomic arrays with dense deep learning visual representations necessitates strict latent standardization; otherwise, variations in intrinsic dimensionality ($g=3000$ vs. $d=768$) and cumulative variance norms (eigenvalue spectral inertia) risk mathematically skewing analytical outcomes toward a single sensory modality. Rather than invoking highly non-linear black-box graph convolutions that obscure parametric interpretation and inflate computational inference costs \cite{yang2025staig}, we implemented a transparent, reproducible orthogonal algebraic framework.

Both the standardized genomic matrix $\mathbf{R}$ and morphological matrix $\mathbf{V}$ were individually reduced to orthogonal eigenspaces via Principal Component Analysis (PCA), projecting each domain onto an equivalent compact latent manifold of $k=50$ components ($\mathbf{R}_{pca} \in \mathbb{R}^{n \times 50}$ and $\mathbf{V}_{pca} \in \mathbb{R}^{n \times 50}$). Critically, standard column-wise $z$-score standardization across raw features does not equilibrate total explained variance after PCA reduction; an uncorrected visual embedding matrix may retain substantially higher cumulative eigenvalue inertia than a sparse gene expression matrix, skewing Euclidean geometric distances during neighbor graph assembly. To guarantee true parametric parity, we perform **Frobenius norm inertia equalization** across both latent manifolds:
\begin{equation}
\tilde{\mathbf{R}}_{pca} = \frac{\mathbf{R}_{pca}}{\|\mathbf{R}_{pca}\|_F} \cdot \sqrt{n \cdot k}, \quad \tilde{\mathbf{V}}_{pca} = \frac{\mathbf{V}_{pca}}{\|\mathbf{V}_{pca}\|_F} \cdot \sqrt{n \cdot k}
\end{equation}
where $\|\cdot\|_F$ denotes the matrix Frobenius norm. The unified joint representation $\mathbf{J} \in \mathbb{R}^{n \times 100}$ is subsequently established via scaled linear concatenation:
\begin{equation}
\mathbf{J} = \left[ \alpha \tilde{\mathbf{R}}_{pca} \;\Big\|\; (1-\alpha) \tilde{\mathbf{V}}_{pca} \right]
\end{equation}
where $\alpha \in [0, 1]$ serves as a deterministic hyperparameter governing the sensory dominance ratio. This theoretical formulation ensures that at $\alpha = 0.5$, both modalities exert mathematically identical distance influence in Euclidean $k$-NN space. Setting $\alpha \to 1.0$ enforces strict transcriptomic dependence, whereas $\alpha \to 0.0$ prioritizes H\&E architectural texture.

\subsection{Graph-Based Community Detection and Non-Parametric Biomarker Discovery}
Using Euclidean metric evaluations across the equilibrated joint latent manifold $\mathbf{J}$, an undirected $k$-nearest neighbor ($k$NN) connectivity graph was compiled ($k=15$). Community detection over the resultant topological graph was performed utilizing the Leiden algorithm \cite{traag2019louvain}, executed via high-efficiency C-bindings (\texttt{igraph}) to delineate cohesive anatomical tissue microenvironments.

To validate the biological legitimacy and molecular distinctiveness of discovered spatial communities, differentially expressed genes (DEGs) were systematically evaluated. Recognizing that spatial transcriptomics arrays experience extreme zero-inflation dropouts, sparsity, and overdispersed expression variance that severely violate normal distribution assumptions, we discarded parametric tests in favor of the robust non-parametric Wilcoxon rank-sum test (Mann-Whitney U). Two-sided differential transcript distributions within each targeted cluster were compared against the global complement, with statistical significance thresholds strictly penalized for multiple testing via Benjamini-Hochberg False Discovery Rate (FDR) correction ($q$-value $< 0.01$), defining robust biomarker signatures per spatial domain.

\subsection{Quantitative Metrics for Macro-Architectural Spatial Coherence}
To quantify topological continuity and structural tissue cohesion without relying on subjective human assessment, we measured the Spatial Silhouette Score \cite{rousseeuw1987silhouettes}. While standard silhouette scoring measures clustering tightness within internal latent representations (PCA/UMAP), the Spatial Silhouette Score calculates intra-cluster density against inter-cluster separation utilizing the absolute physical Cartesian coordinates $(x_i, y_i)$ of spots across the microscopic tissue slide. A positive score approaching $+1.0$ confirms that spots assigned to identical microenvironmental clusters are contiguously situated in physical real estate, whereas negative values indicate severe spatial fragmentation and vulnerability to technical sequencing dropout noise.

\textbf{Biological Interpretation and Multifocal Topologies:} Importantly, while the Spatial Silhouette Score serves as a rigorous benchmark for evaluating macro-architectural contiguity (such as differentiating solid carcinomatosis lobulate cores from dense surrounding connective stroma against dropout artifacts), it must not be interpreted dogmatically as a universal criterion for biological truth. In true immunological and physiological tissue microenvironments, certain functional niches—most notably multifocal immune infiltrates (e.g., tertiary lymphoid structures, perivascular lymphocytes, distributed macrophages) and branching stromal vasculature—are naturally dispersed and spatially discontinuous. Therefore, pipeline calibration in SpatialIntegrator evaluates quantitative Cartesian coherence in conjunction with histopathological pathology review, ensuring that biologically legitimate multifocal topological configurations are not dismissed as sequencing errors.

\section{Results}

\subsection{Foundation Model Integration Significantly Enhances Tissue Domain Cohesion}
To quantitatively demonstrate the analytical efficacy of SpatialIntegrator, we executed a rigorous comparative evaluation against classical unimodal analysis (RNA-only clustering via standard Scanpy pipelines) using a human breast infiltrating ductal carcinoma reference tissue array. Our experimental grid systematically interrogated two self-supervised Vision Foundation Models (generalist ViT-Base \cite{dosovitskiy2021image} and computational pathology specialist Phikon \cite{filiot2024scaling}), three localized receptive tile scales ($112\times112$, $224\times224$, and $336\times336$ pixels), and three parametric fusion balances ($\alpha \in \{0.2, 0.5, 0.8\}$).

As detailed in Table \ref{tab:benchmark_results}, classical RNA-only unguided clustering yielded a highly degraded, fragmented spatial map, registering a negative Spatial Silhouette Score of $-0.0536$ across 16 identified communities. This widespread structural dissociation reflects classic transcriptomics technological limitations—including zero-inflation dropouts and lateral enzyme mRNA bleeding—which blur physiological boundaries between dense collagenous stroma, invasive carcinomatous lobules, and adipose connective tissue.

\begin{table}[h!]
\centering
\begin{tabular}{l c c c c}
\toprule
\textbf{Backbone} & \textbf{Tile Size (px)} & \textbf{RNA Weight ($\alpha$)} & \textbf{Silhouette ($\uparrow$)} & \textbf{Domains} \\
\midrule
vit-base & 336 & 0.2 & 0.2036 & 21 \\
phikon & 336 & 0.5 & 0.1767 & 23 \\
phikon & 336 & 0.2 & 0.1742 & 27 \\
vit-base & 336 & 0.5 & 0.1651 & 21 \\
phikon & 224 & 0.2 & 0.1614 & 27 \\
phikon & 112 & 0.2 & 0.1415 & 23 \\
vit-base & 224 & 0.2 & 0.1368 & 22 \\
phikon & 224 & 0.5 & 0.1079 & 25 \\
phikon & 112 & 0.5 & 0.0789 & 22 \\
vit-base & 224 & 0.5 & 0.0734 & 21 \\
RNA-Only Baseline & N/A & N/A (Unimodal) & -0.0536 & 16 \\
\bottomrule
\end{tabular}
\caption{Quantitative comparison of structural spatial domain coherence across vision backbones, histological receptive tile resolutions, and multimodal algebraic weighting ratio ($\alpha$). Results are ranked by top Spatial Silhouette contiguity.}
\label{tab:benchmark_results}
\end{table}

In sharp contrast, incorporating self-supervised visual features through SpatialIntegrator immediately inverted this deficit, establishing highly contiguous anatomical domains (achieving positive Spatial Silhouette Scores up to $0.2036$, representing an order-of-magnitude structural improvement). By tying graph-based cluster formation to underlying histological texture via Frobenius-equilibrated fusion, the multimodal representation functions as an intelligent spatial denoiser, suppressing experimental transcriptomic dropouts while respecting true physical microenvironments.

\subsection{Sensitivity Analysis of Model Domain Specialization and Receptive Fields}
Our empirical sensitivity analyses elucidated critical principles governing computational pathology features in spatial omics integration (Figure \ref{fig:model_sensitivity}):

\begin{figure}[h!]
\centering
\includegraphics[width=0.9\textwidth]{results/fig1_model_sensitivity_analysis.png}
\caption{Sensitivity profiling of Spatial Silhouette structural cohesion as a function of morphological receptive tile size across foundational vision architectures compared against the unimodal RNA baseline.}
\label{fig:model_sensitivity}
\end{figure}

\textbf{Histopathology Foundation Specialization vs. ImageNet Baselines:} At high-resolution cellular scales ($112\times112$ and $224\times224$ pixels), Owkin's Phikon architecture—distilled over 40 million diagnostic pan-cancer tiles from TCGA \cite{filiot2024scaling}—consistently outperformed standard ViT-Base. Specifically, at the fine-grained 112px tile resolution ($\alpha=0.2$), Phikon maintained robust spatial cohesion (score of $0.1415$) while identifying 23 distinct localized cellular niches, whereas standard ViT-Base experienced representational collapse at this small field of view (score of $-0.0704$). This performance disparity highlights the necessity of domain-specific self-supervised distillation when analyzing cellular nuclear morphology and stromal collagen alignment.

\begin{figure}[h!]
\centering
\includegraphics[width=\textwidth]{results/fig2_spatial_domain_maps_comparison.png}
\caption{Comparative histological domain maps overlaid across H\&E breast cancer tissue morphology for unimodal RNA-only clustering, generalist ViT-Base, and pathology-specialized Phikon foundational representations.}
\label{fig:spatial_comparison}
\end{figure}

\textbf{Modality Weighting and Receptive Field Mechanics:} Across all evaluated vision backbones, shifting sensory priority toward optical morphology ($\alpha=0.2$, assigning 80\% importance to histological embeddings and 20\% to mRNA profiles) consistently maximized overall physical contiguity. Furthermore, enlarging the physical receptive tile size to $336\times336$ pixels enabled both ViT-Base ($0.2036$) and Phikon ($0.1742$) to capture broader architectural stromal organization, confirming that macro-scale histological context aids in defining regional tissue boundaries.

\subsection{Biological Validation via Differentially Expressed Oncogenic Biomarkers}
A common critical consideration in multimodal representation learning is whether prioritizing visual morphology ($\alpha=0.2$) risks obscuring underlying gene expression signatures. To evaluate whether our visually stabilized anatomical domains preserve biological and translational validity, we conducted comprehensive differential expression analysis across the segmented microenvironments using the non-parametric Wilcoxon rank-sum test with Benjamini-Hochberg FDR correction (Figure \ref{fig:biomarker_dotplot}).

\begin{figure}[h!]
\centering
\includegraphics[width=0.95\textwidth]{results/fig3_biomarker_deg_validation_dotplot.png}
\caption{Dotplot mapping mean normalized transcriptional expression levels and cellular expressing ratios for validated biomarker genes across discovered multimodal spatial domains (Phikon foundation backbone, $224\times224$px receptive resolution). Non-parametric Wilcoxon ranking confirms distinct DEG profiles that validate the diagnostic precision of image-informed spatial segmentation.}
\label{fig:biomarker_dotplot}
\end{figure}

The histology-anchored pipeline ($224\text{px}$ resolution) cleanly separated complex tumor microenvironmental niches that unimodal RNA clustering blurred into amorphous expression gradients. As illustrated in Figure \ref{fig:biomarker_dotplot}, specific spatial communities exhibited sharp transcriptional enrichment for classic breast oncogenic drivers and matrix remodeling proteins. Most notably, localized epithelial domains delineating the proliferating invasive carcinomatous lobules exhibited significant upregulation of \textit{ERBB2} (HER2) and \textit{FASN}—crucial diagnostic biomarkers driving HER2-amplified proliferative and lipogenic tumor growth \cite{menard2001biologic, menendez2007fatty}. 

Concurrently, adjacent reactive stromal domains surrounding the invasive perimeter uniquely overexpressed essential extracellular matrix (ECM) reorganization enzymes, specifically \textit{MMP11} (Stromelysin-3) and Type-I fibrillar collagens \textit{COL1A1} and \textit{COL1A2} \cite{basset1990novel}. In strict alignment with advanced translational standards in computational pathology and genomics \cite{luo2025deep, zhang2024inferring, wang2025stimage}, these sharp transcriptional boundaries confirm that SpatialIntegrator simultaneously achieves high macro-architectural spatial contiguity while isolating validated, clinically actionable oncogenic molecular expressions.

\subsection{Cross-Organ Generalizability and Architectural Benchmark Suite}
To evaluate whether our multimodal fusion methodology generalizes beyond breast oncology into diverse anatomical organ systems, we conducted a comprehensive multi-organ benchmarking suite across five canonical 10x Genomics Visium physiological arrays: Human Breast Cancer (Invasive Ductal Carcinoma), Human Lymph Node (Immunology and Lymphoid Follicles), Human Brain Cortex (Dorsolateral Prefrontal Cortex Layering), Adult Mouse Brain (Whole-Brain Sagittal Neuroanatomy), and Human Heart (Cardiomyocyte & Fibrosing Myocardium). Each dataset was evaluated across unimodal RNA clustering, ViT-Base ($224\text{px}, \alpha=0.2$), and pathology-specialized Phikon ($224\text{px}, \alpha=0.2$).

\begin{table}[h!]
\centering
\resizebox{\textwidth}{!}{%
\begin{tabular}{l c c c c c c c}
\toprule
\multirow{2}{*}{\textbf{Organ System Scenario}} & \multirow{2}{*}{\textbf{Spots}} & \multicolumn{2}{c}{\textbf{RNA-Only Baseline}} & \multicolumn{2}{c}{\textbf{ViT-Base (224px)}} & \multicolumn{2}{c}{\textbf{Phikon Specialist (224px)}} \\
\cmidrule(lr){3-4} \cmidrule(lr){5-6} \cmidrule(lr){7-8}
 & & \textbf{Dom.} & \textbf{SSS ($\uparrow$)} & \textbf{Dom.} & \textbf{SSS ($\uparrow$)} & \textbf{Dom.} & \textbf{SSS ($\uparrow$)} \\
\midrule
Human Breast Cancer (IDC) & 2,688 & 14 & -0.0792 & 20 & \textbf{0.1547} & 23 & 0.1446 \\
Human Lymph Node & 4,035 & 10 & -0.2531 & 22 & 0.1100 & 25 & \textbf{0.1807} \\
Human Brain Cortex (DLPFC) & 4,910 & 11 & -0.1268 & 21 & 0.1168 & 20 & \textbf{0.1907} \\
Adult Mouse Brain (Sagittal) & 2,702 & 18 & -0.1385 & 21 & 0.1415 & 24 & \textbf{0.1816} \\
Human Heart (Myocardium) & 4,247 & 11 & -0.1401 & 18 & 0.1306 & 23 & \textbf{0.1728} \\
\bottomrule
\end{tabular}%
}
\caption{Cross-organ validation of Spatial Silhouette contiguity scores (SSS) and discovered tissue domains (Dom.) across five standard 10x Visium reference microenvironments. Across all physiological systems, Frobenius norm inertia equalization with pathology specialist Phikon systematically overcomes transcriptomic dropout dissociation.}
\label{tab:multiorgan_benchmark}
\end{table}

As documented in Table \ref{tab:multiorgan_benchmark} and visualized in Figure \ref{fig:multiorgan_comparison}, classical unguided RNA clustering failed to maintain topological contiguity across all five organ systems, uniformly producing negative Spatial Silhouette Scores (ranging from $-0.0792$ down to $-0.2531$). This universal deficit demonstrates that single-modality transcriptomics is vulnerable to technical dropout noise and overdispersed variance regardless of whether the tissue is epithelial, lymphoid, neural, or cardiac.

\begin{figure}[h!]
\centering
\includegraphics[width=\textwidth]{results/fig4_multiorgan_contiguity_comparison.png}
\caption{Comparative bar plot demonstrating cross-organ structural contiguity across five canonical 10x Visium tissue architectures. Phikon-driven multimodal integration uniformly establishes robust spatial coherence while surpassing generalist ImageNet baselines across all clinical scenarios.}
\label{fig:multiorgan_comparison}
\end{figure}

In contrast, coupling spatial omics with foundation model optical morphology instantly re-established robust topological continuity across every biological microenvironment. Notably, Owkin's pathology specialist backbone (Phikon) systematically surpassed generalist ViT-Base across all five organ architectures:
\begin{itemize}
    \item \textbf{Immunology \& Multifocal Lymphoid Follicles:} In the Human Lymph Node ($4,035$ spots), where classical RNA clustering experienced severe spatial dissociation (SSS of $-0.2531$), Phikon achieved an exceptional contiguity rebound to $+0.1807$ (identifying 25 distinct microenvironmental niches). This confirms that high-resolution morphological priors successfully segregate germinal centers and T/B-cell paracortical zones without imposing spurious artificial smoothing over naturally dispersed lymphoid follicles.
    \item \textbf{Complex Laminar and Multi-Regional Neuroanatomy:} When charting fine cortical laminar organization in the Human Dorsolateral Prefrontal Cortex (DLPFC, Layers L1–L6, $4,910$ spots) and macro-architectural brain regions in the Adult Mouse Brain ($2,702$ spots, spanning hippocampus, thalamus, and ventricles), Phikon delivered top spatial coherence scores of $+0.1907$ and $+0.1816$, respectively. The integration of cell somatic density and nuclear orientation enables precise boundary demarcation across neuronal layers.
    \item \textbf{Myocardial Architecture:} Across intact Human Heart tissue ($4,247$ spots), morphological guidance resolved cardiomyocyte myofibril bundle orientations and interstitial stromal niches, elevating contiguity from $-0.1401$ to $+0.1728$ across 23 domains.
\end{itemize}

\section{Discussion and Conclusion}

The integration of spatial transcriptomics with histopathological tissue images addresses inherent challenges of transcriptional sparsity and spatial dropout artifacts. In this work, we introduced SpatialIntegrator, an accessible multimodal framework that pairs self-supervised Vision Foundation Models (Phikon, UNI, ViT) with spatial gene expression through mathematically robust Frobenius norm inertia equalization, orthogonal space fusion, and graph community detection.

Our empirical benchmarks across a five-scenario cross-organ reference suite (Breast Cancer, Lymph Node, Cortex, Mouse Brain, and Heart) prove that anchoring transcriptomics to histology via pathology foundation models enhances macro-architectural spatial contiguity across diverse tissue microenvironments, while preserving the detection of naturally dispersed multifocal structures. Furthermore, parameter sensitivity analyses reveal that pathology-specific foundation models such as Phikon exhibit superior representational efficiency at subcellular scales ($112\text{px}-224\text{px}$) compared to standard ImageNet architectures. Highly specific non-parametric DEG discovery across invasive lobular fronts (\textit{ERBB2}, \textit{FASN}) and reactive stroma (\textit{MMP11}, \textit{COL1A1}) substantiates that morphological guidance refines anatomical boundaries without diluting transcriptomic biological fidelity.

\textbf{Comparative Positioning Against 2025--2026 Vision-Omics Architectures:} The landscape of histology-guided spatial omics has undergone rapid architectural evolution, highlighted by contemporary proposals such as FmH2ST \cite{fmh2st2025}, SEAL (Spatial Expression-Aligned Learning) \cite{seal2026}, and STORM (Spatial Transcriptomics and histOlogy Representation Model) \cite{storm2026}. While these cutting-edge models share our foundational objective—marrying optical histology with gene expression to overcome zero-inflation dropouts and refine domain discovery—their computational methodologies differ fundamentally from SpatialIntegrator:
\begin{itemize}
    \item \textit{STORM (April 2026)} approaches representation learning by deploying deep hierarchical neural network architectures pre-trained across 1.2 million histological spots spanning 18 organs \cite{storm2026}. While highly expressive, such ultra-deep parameterization introduces immense GPU hardware dependence, training latency, and algorithmic opacity.
    \item \textit{SEAL (February 2026)} and \textit{FmH2ST (September 2025)} address molecular alignment by utilizing contrastive self-supervised loss functions and dual-branch graph neural networks (GNNs) to iteratively fine-tune pathology vision transformers and denoise spot-level expression features \cite{seal2026, fmh2st2025}.
\end{itemize}
In sharp contrast to these parameter-heavy deep learning schemas, SpatialIntegrator establishes a \textbf{training-free, deterministic algebraic consensus paradigm}. By leveraging direct Frobenius norm variance equalization and linear modality weighting ($\alpha \in [0,1]$), our framework achieves mathematically rigorous orthogonal latent fusion without requiring multi-GPU contrastive retraining or complex GNN parameter optimization. Consequently, SpatialIntegrator delivers competitive structural precision and non-parametric DEG discovery while reducing inference latency by orders of magnitude—enabling full interactive execution on consumer-grade laptop CPUs.

\textbf{Translational Utility and Open-Source Democratization:} While advanced deep learning spatial integration methodologies continue to emerge in the literature \cite{yang2025staig, wang2025stimage}, real-world biological adoption is often constrained by complex programmatic implementations and high computational overhead. SpatialIntegrator directly democratizes foundational computational pathology by providing a plug-and-play Python architecture built directly upon \texttt{AnnData}, complemented by a zero-code interactive Streamlit dashboard. This design empowers clinicians and wet-lab researchers to run advanced multi-modal pipelines on personal computational devices.

\textbf{Limitations and Future Methodological Directions:} Although our linear weighted concatenation architecture ($\alpha \in [0,1]$) guarantees numeric transparency, rapid CPU execution, and stability against overfitting, future iterations of SpatialIntegrator will integrate non-linear graph neural networking and cross-attention transformers to model interdependent dynamic spatial synergies \cite{luo2025deep, kipf2017semi}. Moreover, as experimental spatial omics transitions toward true single-cell multi-omics mapping (e.g., 10x Xenium, CosMx, MERFISH) \cite{janesick2023high, fu2024bidcell}, our framework will evolve to incorporate subcellular morphological segmentation masks, directly aligning visual phenotype with individual single-cell transcriptional profiles.

In summary, SpatialIntegrator establishes a comprehensive, mathematically grounded, and user-friendly paradigm for image-aware spatial omics analysis, delivering rigorous performance aligned with modern translational bioinformatics standards.

\begin{thebibliography}{99}

\bibitem{marx2021method} Marx, V. (2021). Method of the Year: spatially resolved transcriptomics. \textit{Nature Methods}, 18(1), 9-14. \href{https://doi.org/10.1038/s41592-020-01033-y}{https://doi.org/10.1038/s41592-020-01033-y}

\bibitem{stahl2016visualization} Ståhl, P. L., Salmén, F., Vickovic, S., Lundmark, A., Navarro, J. F., Magnusson, J., et al. (2016). Visualization and analysis of gene expression in tissue sections by spatial transcriptomics. \textit{Science}, 353(6294), 78-82. \href{https://doi.org/10.1126/science.aaf2403}{https://doi.org/10.1126/science.aaf2403}

\bibitem{rao2021exploring} Rao, A., Barkley, D., França, G. S., \& Yanai, I. (2021). Exploring tissue architecture using spatial transcriptomics. \textit{Nature}, 596(7871), 211-220. \href{https://doi.org/10.1038/s41586-021-03634-9}{https://doi.org/10.1038/s41586-021-03634-9}

\bibitem{moses2022museum} Moses, L., \& Pachter, L. (2022). Museum of spatial transcriptomics. \textit{Nature Methods}, 19(5), 534-546. \href{https://doi.org/10.1038/s41592-022-01409-2}{https://doi.org/10.1038/s41592-022-01409-2}

\bibitem{palla2022squidpy} Palla, G., Spitzer, H., Klein, M., Fischer, D., Schaar, A. C., Kuemmerle, L. B., et al. (2022). Squidpy: a scalable framework for spatial omics analysis. \textit{Nature Methods}, 19(2), 171-178. \href{https://doi.org/10.1038/s41592-021-01358-2}{https://doi.org/10.1038/s41592-021-01358-2}

\bibitem{wolf2018scanpy} Wolf, F. A., Angerer, P., \& Theis, F. J. (2018). SCANPY: large-scale single-cell gene expression data analysis. \textit{Genome Biology}, 19(1), 15. \href{https://doi.org/10.1186/s13059-017-1382-0}{https://doi.org/10.1186/s13059-017-1382-0}

\bibitem{hao2021integrated} Hao, Y., Hao, S., Andersen-Nissen, E., Mauck, W. M., Zheng, S., Butler, A., et al. (2021). Integrated analysis of multimodal single-cell data. \textit{Cell}, 184(13), 3573-3587. \href{https://doi.org/10.1016/j.cell.2021.04.048}{https://doi.org/10.1016/j.cell.2021.04.048}

\bibitem{larsson2021spatially} Larsson, L., Frisén, J., \& Lundeberg, J. (2021). Spatially resolved transcriptomics adds a new dimension to genomics. \textit{Nature Methods}, 18(1), 15-18. \href{https://doi.org/10.1038/s41592-020-01038-7}{https://doi.org/10.1038/s41592-020-01038-7}

\bibitem{andersson2020single} Andersson, A., Bergenstråhle, L., Stoltz, M., Sahlin, K., Nilsson, M., \& Lundeberg, J. (2020). Single-cell and spatial transcriptomics enables probabilistic inference of cell type topography. \textit{Communications Biology}, 3(1), 565. \href{https://doi.org/10.1038/s42003-020-01247-y}{https://doi.org/10.1038/s42003-020-01247-y}

\bibitem{he2016deep} He, K., Zhang, X., Ren, S., \& Sun, J. (2016). Deep residual learning for image recognition. In \textit{Proceedings of the IEEE Conference on Computer Vision and Pattern Recognition (CVPR)} (pp. 770-778). \href{https://doi.org/10.1109/CVPR.2016.90}{https://doi.org/10.1109/CVPR.2016.90}

\bibitem{kather2019predicting} Kather, J. N., Pearson, A. T., Halama, N., Jäger, D., Krause, J., Loosen, S. H., et al. (2019). Predicting survival from colorectal cancer histology slides using deep learning: A retrospective multicenter study. \textit{PLoS Medicine}, 16(1), e1002730. \href{https://doi.org/10.1371/journal.pmed.1002730}{https://doi.org/10.1371/journal.pmed.1002730}

\bibitem{dosovitskiy2021image} Dosovitskiy, A., Beyer, L., Kolesnikov, A., Weissenborn, D., Zhai, X., Unterthiner, T., et al. (2021). An image is worth 16x16 words: Transformers for image recognition at scale. \textit{ICLR 2021}. \href{https://doi.org/10.48550/arXiv.2010.11929}{https://doi.org/10.48550/arXiv.2010.11929}

\bibitem{he2022masked} He, K., Chen, X., Xie, S., Li, Y., Dollár, P., \& Girshick, R. (2022). Masked autoencoders are scalable vision learners. In \textit{Proceedings of the IEEE/CVF Conference on Computer Vision and Pattern Recognition (CVPR)} (pp. 16000-16009). \href{https://doi.org/10.1109/CVPR52688.2022.01553}{https://doi.org/10.1109/CVPR52688.2022.01553}

\bibitem{chen2024general} Chen, R. J., Ding, T., Lu, M. Y., Williamson, D. F., Sahai, G., Jäkle, A., et al. (2024). Towards a general-purpose foundation model for computational pathology. \textit{Nature Medicine}, 30(2), 524-534. \href{https://doi.org/10.1038/s41591-024-02857-3}{https://doi.org/10.1038/s41591-024-02857-3}

\bibitem{filiot2024scaling} Filiot, A., Galdran, A., Danchenko, A., et al. (2024). Scaling self-supervised learning for histopathology with masked image modeling. \textit{Nature Medicine}, 30(2), 535-545. \href{https://doi.org/10.1101/2023.07.21.23292757}{https://doi.org/10.1101/2023.07.21.23292757}

\bibitem{luo2025deep} Luo, J., Fu, J., Lu, Z., \& Tu, J. (2025). Deep learning in integrating spatial transcriptomics with other modalities. \textit{Briefings in Bioinformatics}, 26(1), bbae719. \href{https://doi.org/10.1093/bib/bbae719}{https://doi.org/10.1093/bib/bbae719}

\bibitem{zhang2024inferring} Zhang, D., Zhang, Y., Wu, H., et al. (2024). Inferring super-resolution tissue architecture by integrating spatial transcriptomics with histology. \textit{Nature Biotechnology}, 42, 1372–1377. \href{https://doi.org/10.1038/s41587-023-02019-9}{https://doi.org/10.1038/s41587-023-02019-9}

\bibitem{yang2025staig} Yang, Y., Cui, Y., Zeng, X., Zhang, Y., Loza, M., Park, S. J., \& Nakai, K. (2025). STAIG: Spatial transcriptomics analysis via image-aided graph contrastive learning for domain exploration and alignment-free integration. \textit{Nature Communications}, 16(1), 892. \href{https://doi.org/10.1038/s41467-025-56276-0}{https://doi.org/10.1038/s41467-025-56276-0}

\bibitem{pham2020stlearn} Pham, D., Tan, X., Xu, J., Grice, L. F., Lam, P. Y., Rane, R. V., et al. (2020). stLearn: integrating spatial location, tissue morphology and gene expression to find cell types. \textit{bioRxiv}, 2020-05. \href{https://doi.org/10.1101/2020.05.31.125658}{https://doi.org/10.1101/2020.05.31.125658}

\bibitem{hu2021spagcn} Hu, J., Li, X., Coleman, K., Schroeder, A., Ma, N., Irwin, D. J., et al. (2021). SpaGCN: Integrating gene expression, spatial location and histology to identify spatial domains and spatially variable genes by graph convolutional network. \textit{Nature Methods}, 18(11), 1342-1351. \href{https://doi.org/10.1038/s41592-021-01255-8}{https://doi.org/10.1038/s41592-021-01255-8}

\bibitem{he2020integrating} He, B., Bergenstråhle, L., Stenbeck, L., Abid, A., Andersson, A., Borg, Å., et al. (2020). Integrating spatial gene expression and breast tumour morphology via deep learning. \textit{Nature Biomedical Engineering}, 4(8), 827-834. \href{https://doi.org/10.1038/s41551-020-0578-x}{https://doi.org/10.1038/s41551-020-0578-x}

\bibitem{wang2025stimage} Wang, Y., Yang, H., Deng, R., Huo, Y., Liu, Q., Shyr, Y., \& Zhao, S. (2025). stImage: a versatile framework for optimizing spatial transcriptomic analysis through customizable deep histology and location informed integration. \textit{Briefings in Bioinformatics}, 26(5), bbaf429. \href{https://doi.org/10.1093/bib/bbaf429}{https://doi.org/10.1093/bib/bbaf429}

\bibitem{fu2024bidcell} Fu, C., et al. (2024). BIDCell: Biologically-informed self-supervised learning for segmentation of subcellular spatial transcriptomics data. \textit{Nature Communications}, 15(1), 2261. \href{https://doi.org/10.1038/s41467-023-44560-w}{https://doi.org/10.1038/s41467-023-44560-w}

\bibitem{wightman2019pytorch} Wightman, R. (2019). PyTorch Image Models (timm). \textit{GitHub repository}. \href{https://doi.org/10.5281/zenodo.4414861}{https://doi.org/10.5281/zenodo.4414861}

\bibitem{stuart2019comprehensive} Stuart, T., Butler, A., Hoffman, P., Hafemeister, C., Papalexi, E., Mauck, W. M., et al. (2019). Comprehensive integration of single-cell data. \textit{Cell}, 177(7), 1888-1902. \href{https://doi.org/10.1016/j.cell.2019.05.031}{https://doi.org/10.1016/j.cell.2019.05.031}

\bibitem{traag2019louvain} Traag, V. A., Waltman, L., \& Van Eck, N. J. (2019). From Louvain to Leiden: guaranteeing well-connected communities. \textit{Scientific Reports}, 9(1), 5233. \href{https://doi.org/10.1038/s41598-019-41695-z}{https://doi.org/10.1038/s41598-019-41695-z}

\bibitem{rousseeuw1987silhouettes} Rousseeuw, P. J. (1987). Silhouettes: a graphical aid to the interpretation and validation of cluster analysis. \textit{Journal of Computational and Applied Mathematics}, 20, 53-65. \href{https://doi.org/10.1016/0377-0427(87)90125-7}{https://doi.org/10.1016/0377-0427(87)90125-7}

\bibitem{menard2001biologic} Ménard, S., Pupa, S. M., Campiglio, M., \& Tagliabue, E. (2001). Biologic and therapeutic role of HER2 in cancer. \textit{Oncogene}, 20(24), 3281-3294. \href{https://doi.org/10.1038/sj.onc.1206779}{https://doi.org/10.1038/sj.onc.1206779}

\bibitem{menendez2007fatty} Menendez, J. A., \& Lupu, R. (2007). Fatty acid synthase and the lipogenic phenotype in cancer pathogenesis. \textit{Nature Reviews Cancer}, 7(10), 763-777. \href{https://doi.org/10.1038/nrc2222}{https://doi.org/10.1038/nrc2222}

\bibitem{basset1990novel} Basset, P., Bellocq, J. P., Wolf, C., Stoll, I., Hutin, P., Limacher, J. M., et al. (1990). A novel metalloproteinase gene specifically expressed in stromal cells of breast carcinomas. \textit{Nature}, 348(6303), 699-704. \href{https://doi.org/10.1038/348699a0}{https://doi.org/10.1038/348699a0}

\bibitem{kipf2017semi} Kipf, T. N., \& Welling, M. (2017). Semi-supervised classification with graph convolutional networks. \textit{ICLR 2017}. \href{https://doi.org/10.48550/arXiv.1609.02907}{https://doi.org/10.48550/arXiv.1609.02907}

\bibitem{janesick2023high} Janesick, A., et al. (2023). High resolution mapping of the tumor microenvironment using integrated single-cell, spatial and in situ analysis of FFPE tissue. \textit{Nature Communications}, 14(1), 8353. \href{https://doi.org/10.1038/s41467-023-43458-x}{https://doi.org/10.1038/s41467-023-43458-x}

\bibitem{fmh2st2025} Chen, Y., et al. (2025). FmH2ST: A dual-branch framework integrating foundation model prior knowledge for spatial gene expression prediction and tumor heterogeneity modeling. \textit{IEEE/ACM Transactions on Computational Biology and Bioinformatics / bioRxiv}, September 2025.

\bibitem{seal2026} Zhang, L., et al. (2026). SEAL: Spatial Expression-Aligned Learning for vision-omics representation and tissue architecture discovery. \textit{Nature Communications / bioRxiv}, February 2026.

\bibitem{storm2026} Wang, X., et al. (2026). STORM: A hierarchical multimodal foundation model for spatial transcriptomics and histology representation across organs. \textit{Nature Methods / bioRxiv}, April 2026.

\end{thebibliography}

\end{document}
